\section{Frequent Losers and Outcomes in Tenders}

\section{The outcomes observed in public tenders with participation by frequent losers can identify evidence of explicit collusion, potentially serving as a screening mechanism for the existence of these practices. As developed in Section 2, an essential principle of effective screening models is their ability to distinguish collusive behaviors from those that could arise from different competitive models.}

\section{This section investigates the properties of this collusion marker and observes the relationship between the presence of frequent losers and the performance of bidding contests in several dimensions, such as prices, the number of participants, and the number of bids. Based on these results, we evaluate the effectiveness of this marker in identifying collusive behavior.}

\section{Possible differences in bidding results are estimated by comparing public tenders with at least one frequent loser and tenders without the presence of a frequent loser. This estimation of differences in outcomes y between purchases i, with and without the participation of a frequent loser, of an item g and in year t has the following baseline specification:}

\vspace{1\baselineskip}
\begin{equation}
\mathbf{y}_{\mathbf{igt}}\mathbf{ = }\mathbf{\beta }\mathbf{losers}_{\mathbf{igt}}\mathbf{+}\mathbf{\alpha }_{\mathbf{g}}\mathbf{+}\mathbf{\lambda }_{\mathbf{t}}\mathbf{+}\mathbf{x\delta }\mathbf{+}\mathbf{\epsilon }_{\mathbf{igt,}}\mathbf{     }\left(\mathbf{2}\right)
\end{equation}


\vspace{1\baselineskip}
\section{where and are fixed effects of purchased items and year dummies, respectively. The variable has a value equal to 1 if the tender has at least one frequent loser and 0 otherwise. Additionally, are the control variables, such as dummies for purchasing units and tender mode.}

\section{The data used come from public purchases made at BEC from January 2009 to December 2019. However, for the purposes of this section, the subsample used in the estimates consists only of items with at least one purchasing process with the presence of frequent losers. The estimates for the negotiated prices are presented in Table 2.}

\begin{table}[H]
\begin{adjustbox}{max width=\textwidth}
\begin{tabular}{p{3.92cm}p{2.57cm}p{2.5cm}p{2.5cm}p{2.5cm}}
\hline
\multicolumn{1}{|p{5.3cm}}{} &
\multicolumn{1}{|p{3.14cm}}{\centering
 (1)} &
\multicolumn{1}{|p{3.02cm}}{\centering
 (2)} &
\multicolumn{1}{|p{3.02cm}}{\centering
 (3)} &
\multicolumn{1}{|p{3.02cm}|}{\centering
 (4)} \\
\hline
\multicolumn{1}{|p{5.3cm}}{} &
\multicolumn{1}{|p{3.14cm}}{\centering
    General} &
\multicolumn{1}{|p{3.02cm}}{\centering
    General} &
\multicolumn{1}{|p{3.02cm}}{\centering
 Pregao} &
\multicolumn{1}{|p{3.02cm}|}{\centering
    Convite} \\
\hline
\multicolumn{1}{|p{5.3cm}}{ losers} &
\multicolumn{1}{|p{3.14cm}}{\centering
.1255$\ast$$\ast$$\ast$} &
\multicolumn{1}{|p{3.02cm}}{\centering
.1011$\ast$$\ast$$\ast$} &
\multicolumn{1}{|p{3.02cm}}{\centering
.1344$\ast$$\ast$$\ast$} &
\multicolumn{1}{|p{3.02cm}|}{\centering
.0315$\ast$$\ast$$\ast$} \\
\hline
\multicolumn{1}{|p{5.3cm}}{} &
\multicolumn{1}{|p{3.14cm}}{\centering
(.014)} &
\multicolumn{1}{|p{3.02cm}}{\centering
(.0122)} &
\multicolumn{1}{|p{3.02cm}}{\centering
(.0184)} &
\multicolumn{1}{|p{3.02cm}|}{\centering
(.0046)} \\
\hline
\multicolumn{1}{|p{5.3cm}}{ convite} &
\multicolumn{1}{|p{3.14cm}}{\centering
-.0449$\ast$$\ast$$\ast$} &
\multicolumn{1}{|p{3.02cm}}{\centering
-.0442$\ast$$\ast$$\ast$} &
\multicolumn{1}{|p{3.02cm}}{} &
\multicolumn{1}{|p{3.02cm}|}{} \\
\hline
\multicolumn{1}{|p{5.3cm}}{} &
\multicolumn{1}{|p{3.14cm}}{\centering
(.0057)} &
\multicolumn{1}{|p{3.02cm}}{\centering
(.0056)} &
\multicolumn{1}{|p{3.02cm}}{} &
\multicolumn{1}{|p{3.02cm}|}{} \\
\hline
\multicolumn{1}{|p{5.3cm}}{ pregao} &
\multicolumn{1}{|p{3.14cm}}{\centering
.4865$\ast$$\ast$$\ast$} &
\multicolumn{1}{|p{3.02cm}}{\centering
.5653$\ast$$\ast$$\ast$} &
\multicolumn{1}{|p{3.02cm}}{} &
\multicolumn{1}{|p{3.02cm}|}{} \\
\hline
\multicolumn{1}{|p{5.3cm}}{} &
\multicolumn{1}{|p{3.14cm}}{\centering
(.0177)} &
\multicolumn{1}{|p{3.02cm}}{\centering
(.0208)} &
\multicolumn{1}{|p{3.02cm}}{} &
\multicolumn{1}{|p{3.02cm}|}{} \\
\hline
\multicolumn{1}{|p{5.3cm}}{ lquantity} &
\multicolumn{1}{|p{3.14cm}}{\centering
-.2666$\ast$$\ast$$\ast$} &
\multicolumn{1}{|p{3.02cm}}{\centering
-.2988$\ast$$\ast$$\ast$} &
\multicolumn{1}{|p{3.02cm}}{\centering
-.4781$\ast$$\ast$$\ast$} &
\multicolumn{1}{|p{3.02cm}|}{\centering
-.2118$\ast$$\ast$$\ast$} \\
\hline
\multicolumn{1}{|p{5.3cm}}{} &
\multicolumn{1}{|p{3.14cm}}{\centering
(.0098)} &
\multicolumn{1}{|p{3.02cm}}{\centering
(.0098)} &
\multicolumn{1}{|p{3.02cm}}{\centering
(.0175)} &
\multicolumn{1}{|p{3.02cm}|}{\centering
(.0109)} \\
\hline
\multicolumn{1}{|p{5.3cm}}{ \_cons} &
\multicolumn{1}{|p{3.14cm}}{\centering
3.6905$\ast$$\ast$$\ast$} &
\multicolumn{1}{|p{3.02cm}}{\centering
4.0337$\ast$$\ast$$\ast$} &
\multicolumn{1}{|p{3.02cm}}{\centering
5.5803$\ast$$\ast$$\ast$} &
\multicolumn{1}{|p{3.02cm}|}{\centering
3.7513$\ast$$\ast$$\ast$} \\
\hline
\multicolumn{1}{|p{5.3cm}}{} &
\multicolumn{1}{|p{3.14cm}}{\centering
(.0352)} &
\multicolumn{1}{|p{3.02cm}}{\centering
(.0703)} &
\multicolumn{1}{|p{3.02cm}}{\centering
(.1422)} &
\multicolumn{1}{|p{3.02cm}|}{\centering
(.0825)} \\
\hline
\multicolumn{1}{|p{5.3cm}}{ Observations} &
\multicolumn{1}{|p{3.14cm}}{\centering
1671773} &
\multicolumn{1}{|p{3.02cm}}{\centering
1671773} &
\multicolumn{1}{|p{3.02cm}}{\centering
474219} &
\multicolumn{1}{|p{3.02cm}|}{\centering
924725} \\
\hline
\multicolumn{1}{|p{5.3cm}}{ R-squared} &
\multicolumn{1}{|p{3.14cm}}{\centering
.2085} &
\multicolumn{1}{|p{3.02cm}}{\centering
.2596} &
\multicolumn{1}{|p{3.02cm}}{\centering
.4219} &
\multicolumn{1}{|p{3.02cm}|}{\centering
.2066} \\
\hline
\multicolumn{1}{|p{5.3cm}}{Item Dummies} &
\multicolumn{1}{|p{3.14cm}}{\centering
YES} &
\multicolumn{1}{|p{3.02cm}}{\centering
YES} &
\multicolumn{1}{|p{3.02cm}}{\centering
YES} &
\multicolumn{1}{|p{3.02cm}|}{\centering
YES} \\
\hline
\multicolumn{1}{|p{5.3cm}}{Year Dummies} &
\multicolumn{1}{|p{3.14cm}}{\centering
YES} &
\multicolumn{1}{|p{3.02cm}}{\centering
YES} &
\multicolumn{1}{|p{3.02cm}}{\centering
YES} &
\multicolumn{1}{|p{3.02cm}|}{\centering
YES} \\
\hline
\multicolumn{1}{|p{5.3cm}}{PBU Dummies} &
\multicolumn{1}{|p{3.14cm}}{\centering
NO} &
\multicolumn{1}{|p{3.02cm}}{\centering
YES} &
\multicolumn{1}{|p{3.02cm}}{\centering
YES} &
\multicolumn{1}{|p{3.02cm}|}{\centering
YES} \\
\hline
\multicolumn{5}{|p{17.49cm}|}{\textit{Standard errors are in parentheses}} \\
\hline
\multicolumn{5}{|p{17.49cm}|}{\textit{$\ast$$\ast$$\ast$ p<.01, $\ast$$\ast$ p<.05, $\ast$ p<.1}} \\
\hline
\end{tabular}
\end{adjustbox}
\caption{With vs. Without Frequent Losers}
\label{tab:vs_frequent_losers}\end{table}
\vspace{17\baselineskip}
\section{It can be observed that the negotiated prices are consistently higher in tenders with the presence of frequent losers. Considering variations in the baseline specification containing all the bidding modes, the negotiated prices, on average, are between 10.64% and 13.37% higher in bids with losers.}

\section{The results also suggest that prices in pregoes tend to be more affected than prices in convites by the participation of frequent losers. Considering only pregoes, it is possible to observe that prices tend to be 14.39% higher in bids that include losers than when losers are not present; for convites, prices are only 3.20% higher on the same comparative basis.}

\section{Possible explanations for this difference might be related to (i) the degree of discretion of the public official in charge of the tender and (ii) the possibility for suppliers to affect the dynamics and outcome of the process.}

\section{In convites, the process consists of merely opening (virtual) envelopes containing participating suppliers' proposals. After making the bids public, the auctioneer's role is restricted to declaring the firm that offered the lowest price the winner. Thus, the interaction between the participating firms and public officers is minimal during the purchase procedure.}

\section{In the case of the pregao, the public officer in charge has a more active role in the process, and there is more room for firms to influence the process. In addition to revealing the initial bids in the first phase, the auctioneer is responsible for coordinating and monitoring firms' bids during the auction phase.}

\section{During this auction phase, firms and PBUs interact through the real-time bidding process and real-time chat. A great variety of information is exchanged through this communication channel, such as confirming product specifications, complaints about competitors' performance, and even dissatisfaction with the way the contest is conducted.}

\section{At the end of the auction phase, there is another moment of interaction between suppliers and PBUs. There is an ex-post negotiation phase in which public and private parties can direct bargaining for the lowest price. Thus, in pregao, there is more room for manipulation. This situation may lead to higher prices in this type of tender.}

\section{Higher prices in tenders with frequent losers may suggest that there is less competitiveness or aggressiveness for suppliers. This fact cannot be verified in the estimates of Tables 3 and 4. In these tables, we compare tenders with and without frequent losers in terms of the number of firms and bids, respectively.}

\vspace{1\baselineskip}
\begin{center}
Table 3. Number of Firms: With vs. Without Frequent Losers
\end{center}


\begin{table}[H]
\begin{adjustbox}{max width=\textwidth}
\begin{tabular}{p{3.98cm}p{2.52cm}p{2.53cm}p{2.53cm}p{2.45cm}}
\hline
\multicolumn{1}{|p{5.3cm}}{} &
\multicolumn{1}{|p{3.14cm}}{\centering
 (1)} &
\multicolumn{1}{|p{3.02cm}}{\centering
 (2)} &
\multicolumn{1}{|p{3.02cm}}{\centering
 (3)} &
\multicolumn{1}{|p{3.02cm}|}{\centering
 (4)} \\
\hline
\multicolumn{1}{|p{5.3cm}}{} &
\multicolumn{1}{|p{3.14cm}}{\centering
    General} &
\multicolumn{1}{|p{3.02cm}}{\centering
    General} &
\multicolumn{1}{|p{3.02cm}}{\centering
 Pregao} &
\multicolumn{1}{|p{3.02cm}|}{\centering
    Convite} \\
\hline
\multicolumn{1}{|p{5.3cm}}{ losers} &
\multicolumn{1}{|p{3.14cm}}{\centering
.329$\ast$$\ast$$\ast$} &
\multicolumn{1}{|p{3.02cm}}{\centering
.3256$\ast$$\ast$$\ast$} &
\multicolumn{1}{|p{3.02cm}}{\centering
.3104$\ast$$\ast$$\ast$} &
\multicolumn{1}{|p{3.02cm}|}{\centering
.2926$\ast$$\ast$$\ast$} \\
\hline
\multicolumn{1}{|p{5.3cm}}{} &
\multicolumn{1}{|p{3.14cm}}{\centering
(.0052)} &
\multicolumn{1}{|p{3.02cm}}{\centering
(.0048)} &
\multicolumn{1}{|p{3.02cm}}{\centering
(.008)} &
\multicolumn{1}{|p{3.02cm}|}{\centering
(.0041)} \\
\hline
\multicolumn{1}{|p{5.3cm}}{ convite} &
\multicolumn{1}{|p{3.14cm}}{\centering
-.7742$\ast$$\ast$$\ast$} &
\multicolumn{1}{|p{3.02cm}}{\centering
-.7755$\ast$$\ast$$\ast$} &
\multicolumn{1}{|p{3.02cm}}{} &
\multicolumn{1}{|p{3.02cm}|}{} \\
\hline
\multicolumn{1}{|p{5.3cm}}{} &
\multicolumn{1}{|p{3.14cm}}{\centering
(.0069)} &
\multicolumn{1}{|p{3.02cm}}{\centering
(.007)} &
\multicolumn{1}{|p{3.02cm}}{} &
\multicolumn{1}{|p{3.02cm}|}{} \\
\hline
\multicolumn{1}{|p{5.3cm}}{ pregao} &
\multicolumn{1}{|p{3.14cm}}{\centering
-.1181$\ast$$\ast$$\ast$} &
\multicolumn{1}{|p{3.02cm}}{\centering
-.1228$\ast$$\ast$$\ast$} &
\multicolumn{1}{|p{3.02cm}}{} &
\multicolumn{1}{|p{3.02cm}|}{} \\
\hline
\multicolumn{1}{|p{5.3cm}}{} &
\multicolumn{1}{|p{3.14cm}}{\centering
(.0081)} &
\multicolumn{1}{|p{3.02cm}}{\centering
(.0099)} &
\multicolumn{1}{|p{3.02cm}}{} &
\multicolumn{1}{|p{3.02cm}|}{} \\
\hline
\multicolumn{1}{|p{5.3cm}}{ lquantity} &
\multicolumn{1}{|p{3.14cm}}{\centering
.1504$\ast$$\ast$$\ast$} &
\multicolumn{1}{|p{3.02cm}}{\centering
.151$\ast$$\ast$$\ast$} &
\multicolumn{1}{|p{3.02cm}}{\centering
.0925$\ast$$\ast$$\ast$} &
\multicolumn{1}{|p{3.02cm}|}{\centering
.1817$\ast$$\ast$$\ast$} \\
\hline
\multicolumn{1}{|p{5.3cm}}{} &
\multicolumn{1}{|p{3.14cm}}{\centering
(.0021)} &
\multicolumn{1}{|p{3.02cm}}{\centering
(.0022)} &
\multicolumn{1}{|p{3.02cm}}{\centering
(.0028)} &
\multicolumn{1}{|p{3.02cm}|}{\centering
(.0028)} \\
\hline
\multicolumn{1}{|p{5.3cm}}{ \_cons} &
\multicolumn{1}{|p{3.14cm}}{\centering
.9159$\ast$$\ast$$\ast$} &
\multicolumn{1}{|p{3.02cm}}{\centering
1.2013$\ast$$\ast$$\ast$} &
\multicolumn{1}{|p{3.02cm}}{\centering
1.0851$\ast$$\ast$$\ast$} &
\multicolumn{1}{|p{3.02cm}|}{\centering
.8704$\ast$$\ast$$\ast$} \\
\hline
\multicolumn{1}{|p{5.3cm}}{} &
\multicolumn{1}{|p{3.14cm}}{\centering
(.0093)} &
\multicolumn{1}{|p{3.02cm}}{\centering
(.028)} &
\multicolumn{1}{|p{3.02cm}}{\centering
(.0451)} &
\multicolumn{1}{|p{3.02cm}|}{\centering
(.0624)} \\
\hline
\multicolumn{1}{|p{5.3cm}}{ Observations} &
\multicolumn{1}{|p{3.14cm}}{\centering
1670719} &
\multicolumn{1}{|p{3.02cm}}{\centering
1670719} &
\multicolumn{1}{|p{3.02cm}}{\centering
473819} &
\multicolumn{1}{|p{3.02cm}|}{\centering
924336} \\
\hline
\multicolumn{1}{|p{5.3cm}}{ R-squared} &
\multicolumn{1}{|p{3.14cm}}{\centering
.3675} &
\multicolumn{1}{|p{3.02cm}}{\centering
.3882} &
\multicolumn{1}{|p{3.02cm}}{\centering
.2176} &
\multicolumn{1}{|p{3.02cm}|}{\centering
.2699} \\
\hline
\multicolumn{1}{|p{5.3cm}}{Items Dummies} &
\multicolumn{1}{|p{3.14cm}}{\centering
YES} &
\multicolumn{1}{|p{3.02cm}}{\centering
YES} &
\multicolumn{1}{|p{3.02cm}}{\centering
YES} &
\multicolumn{1}{|p{3.02cm}|}{\centering
YES} \\
\hline
\multicolumn{1}{|p{5.3cm}}{Year Dummies} &
\multicolumn{1}{|p{3.14cm}}{\centering
YES} &
\multicolumn{1}{|p{3.02cm}}{\centering
YES} &
\multicolumn{1}{|p{3.02cm}}{\centering
YES} &
\multicolumn{1}{|p{3.02cm}|}{\centering
YES} \\
\hline
\multicolumn{1}{|p{5.3cm}}{PBU Dummies} &
\multicolumn{1}{|p{3.14cm}}{\centering
NO} &
\multicolumn{1}{|p{3.02cm}}{\centering
YES} &
\multicolumn{1}{|p{3.02cm}}{\centering
YES} &
\multicolumn{1}{|p{3.02cm}|}{\centering
YES} \\
\hline
\multicolumn{5}{|p{17.49cm}|}{\textit{Standard errors are in parentheses}} \\
\hline
\multicolumn{5}{|p{17.49cm}|}{\textit{$\ast$$\ast$$\ast$ p<.01, $\ast$$\ast$ p<.05, $\ast$ p<.1}} \\
\hline
\end{tabular}
\end{adjustbox}
\end{table}
\vspace{18\baselineskip}
\section{More firms participate and more bids are offered in tenders with frequent losers than in bids without their presence. Between 33.99% and 38.96% more firms participate in tenders with losers, with an average of 25.51% to 34.35% more bids offered.}

\vspace{2\baselineskip}
\begin{center}
Table 4. Number of Bids: With vs. Without Frequent Losers
\end{center}


\begin{table}[H]
\begin{adjustbox}{max width=\textwidth}
\begin{tabular}{p{4.21cm}p{2.5cm}p{2.4cm}p{2.4cm}p{2.5cm}}
\hline
\multicolumn{1}{|p{4.21cm}}{} &
\multicolumn{1}{|p{2.5cm}}{\centering
 (1)} &
\multicolumn{1}{|p{2.4cm}}{\centering
 (2)} &
\multicolumn{1}{|p{2.4cm}}{\centering
 (3)} &
\multicolumn{1}{|p{2.5cm}|}{\centering
 (4)} \\
\hline
\multicolumn{1}{|p{4.21cm}}{} &
\multicolumn{1}{|p{2.5cm}}{\centering
    General} &
\multicolumn{1}{|p{2.4cm}}{\centering
    General} &
\multicolumn{1}{|p{2.4cm}}{\centering
 Pregao} &
\multicolumn{1}{|p{2.5cm}|}{\centering
    Convite} \\
\hline
\multicolumn{1}{|p{4.21cm}}{ losers} &
\multicolumn{1}{|p{2.5cm}}{\centering
.2953$\ast$$\ast$$\ast$} &
\multicolumn{1}{|p{2.4cm}}{\centering
.2937$\ast$$\ast$$\ast$} &
\multicolumn{1}{|p{2.4cm}}{\centering
.2272$\ast$$\ast$$\ast$} &
\multicolumn{1}{|p{2.5cm}|}{\centering
.2926$\ast$$\ast$$\ast$} \\
\hline
\multicolumn{1}{|p{4.21cm}}{} &
\multicolumn{1}{|p{2.5cm}}{\centering
(.0062)} &
\multicolumn{1}{|p{2.4cm}}{\centering
(.006)} &
\multicolumn{1}{|p{2.4cm}}{\centering
(.0108)} &
\multicolumn{1}{|p{2.5cm}|}{\centering
(.0041)} \\
\hline
\multicolumn{1}{|p{4.21cm}}{ convite} &
\multicolumn{1}{|p{2.5cm}}{\centering
.0314$\ast$$\ast$$\ast$} &
\multicolumn{1}{|p{2.4cm}}{\centering
.0249$\ast$$\ast$$\ast$} &
\multicolumn{1}{|p{2.4cm}}{} &
\multicolumn{1}{|p{2.5cm}|}{} \\
\hline
\multicolumn{1}{|p{4.21cm}}{} &
\multicolumn{1}{|p{2.5cm}}{\centering
(.0093)} &
\multicolumn{1}{|p{2.4cm}}{\centering
(.0093)} &
\multicolumn{1}{|p{2.4cm}}{} &
\multicolumn{1}{|p{2.5cm}|}{} \\
\hline
\multicolumn{1}{|p{4.21cm}}{ pregao} &
\multicolumn{1}{|p{2.5cm}}{\centering
.931$\ast$$\ast$$\ast$} &
\multicolumn{1}{|p{2.4cm}}{\centering
.9225$\ast$$\ast$$\ast$} &
\multicolumn{1}{|p{2.4cm}}{} &
\multicolumn{1}{|p{2.5cm}|}{} \\
\hline
\multicolumn{1}{|p{4.21cm}}{} &
\multicolumn{1}{|p{2.5cm}}{\centering
(.01)} &
\multicolumn{1}{|p{2.4cm}}{\centering
(.0126)} &
\multicolumn{1}{|p{2.4cm}}{} &
\multicolumn{1}{|p{2.5cm}|}{} \\
\hline
\multicolumn{1}{|p{4.21cm}}{ lquantity} &
\multicolumn{1}{|p{2.5cm}}{\centering
.169$\ast$$\ast$$\ast$} &
\multicolumn{1}{|p{2.4cm}}{\centering
.1674$\ast$$\ast$$\ast$} &
\multicolumn{1}{|p{2.4cm}}{\centering
.1059$\ast$$\ast$$\ast$} &
\multicolumn{1}{|p{2.5cm}|}{\centering
.1817$\ast$$\ast$$\ast$} \\
\hline
\multicolumn{1}{|p{4.21cm}}{} &
\multicolumn{1}{|p{2.5cm}}{\centering
(.0026)} &
\multicolumn{1}{|p{2.4cm}}{\centering
(.0027)} &
\multicolumn{1}{|p{2.4cm}}{\centering
(.004)} &
\multicolumn{1}{|p{2.5cm}|}{\centering
(.0028)} \\
\hline
\multicolumn{1}{|p{4.21cm}}{ \_cons} &
\multicolumn{1}{|p{2.5cm}}{\centering
.7849$\ast$$\ast$$\ast$} &
\multicolumn{1}{|p{2.4cm}}{\centering
1.5003$\ast$$\ast$$\ast$} &
\multicolumn{1}{|p{2.4cm}}{\centering
2.17$\ast$$\ast$$\ast$} &
\multicolumn{1}{|p{2.5cm}|}{\centering
.8704$\ast$$\ast$$\ast$} \\
\hline
\multicolumn{1}{|p{4.21cm}}{} &
\multicolumn{1}{|p{2.5cm}}{\centering
(.0114)} &
\multicolumn{1}{|p{2.4cm}}{\centering
(.0478)} &
\multicolumn{1}{|p{2.4cm}}{\centering
(.077)} &
\multicolumn{1}{|p{2.5cm}|}{\centering
(.0624)} \\
\hline
\multicolumn{1}{|p{4.21cm}}{ Observations} &
\multicolumn{1}{|p{2.5cm}}{\centering
1670719} &
\multicolumn{1}{|p{2.4cm}}{\centering
1670719} &
\multicolumn{1}{|p{2.4cm}}{\centering
473819} &
\multicolumn{1}{|p{2.5cm}|}{\centering
924336} \\
\hline
\multicolumn{1}{|p{4.21cm}}{ R-squared} &
\multicolumn{1}{|p{2.5cm}}{\centering
.2945} &
\multicolumn{1}{|p{2.4cm}}{\centering
.3117} &
\multicolumn{1}{|p{2.4cm}}{\centering
.1351} &
\multicolumn{1}{|p{2.5cm}|}{\centering
.2698} \\
\hline
\multicolumn{1}{|p{4.21cm}}{Items Dummies} &
\multicolumn{1}{|p{2.5cm}}{\centering
YES} &
\multicolumn{1}{|p{2.4cm}}{\centering
YES} &
\multicolumn{1}{|p{2.4cm}}{\centering
YES} &
\multicolumn{1}{|p{2.5cm}|}{\centering
YES} \\
\hline
\multicolumn{1}{|p{4.21cm}}{Year Dummies} &
\multicolumn{1}{|p{2.5cm}}{\centering
YES} &
\multicolumn{1}{|p{2.4cm}}{\centering
YES} &
\multicolumn{1}{|p{2.4cm}}{\centering
YES} &
\multicolumn{1}{|p{2.5cm}|}{\centering
YES} \\
\hline
\multicolumn{1}{|p{4.21cm}}{PBU Dummies} &
\multicolumn{1}{|p{2.5cm}}{\centering
NO} &
\multicolumn{1}{|p{2.4cm}}{\centering
YES} &
\multicolumn{1}{|p{2.4cm}}{\centering
YES} &
\multicolumn{1}{|p{2.5cm}|}{\centering
YES} \\
\hline
\multicolumn{5}{|p{14.01cm}|}{\textit{Standard errors are in parentheses}} \\
\hline
\multicolumn{5}{|p{14.01cm}|}{\textit{$\ast$$\ast$$\ast$ p<.01, $\ast$$\ast$ p<.05, $\ast$ p<.1}} \\
\hline
\end{tabular}
\end{adjustbox}
\end{table}
\vspace{19\baselineskip}
\section{These results are compatible with the literature on cartels and other illicit purchases. As ) notes, cartel participants' behavior also responds strategically to competition policy enforcement. A higher number of participating firms and bidders could be associated with attempting to build a mechanism to conceal possible illicit competitive practices.}

\section{It is important to note that, assuming there is no correlation between the companies' cost structure and the presence of frequent losers, the simultaneous occurrence of higher prices and a higher number of bids and participants in public tenders is inconsistent with expectations for a model of competition or tacit collusion.}

\section{This scenario might mean that the screening proposal based on frequent losers presents two essential properties. First, it shows that it is capable of distinguishing collusive behaviors from those that would be expected in competition. Second and more importantly, because the observed behavior reveals the intention to hide the collusion, it is a marker that identifies explicit and non-tacit collusion, a recurrent problem in screening models ).}

\section{To establish the reliability of this screening model, however, it is necessary to highlight possible alternative explanations for the presence of frequent losers in public procurement. A first alternative hypothesis is the existence of a learning curve for companies to become victorious at BEC. Supplying goods and services to the state would require suppliers to develop specific processes and legal adjustments that would demand consistent efforts over some time.}

\section{If there were such a learning curve, then companies that eventually win the first bid would spend a considerable period of time losing until they become victorious. This situation may be the case for some suppliers; however, the data suggest that this is not the general case. Figure 9 shows the number of contests up to a company's first victory, limited to 100 contests for better viewing. Almost 30% of companies win one of the first three events in which they participate.}

\vspace{1\baselineskip}
\begin{figure}[H]
\includegraphics[width=10.82cm,height=7.3cm]{./images/image9.png}
\caption{-Number of Participations in Tenders until the First Victory}
\label{fig:number_participations_tenders_first_victory}
\end{figure}


\vspace{1\baselineskip}
\section{In addition, as explained in Section 3, it is known that a small number of companies win public tenders. Thus, a learning curve hypothesis seems unlikely, except for a few companies that achieve very frequent wins in bids.}

\section{There is an additional alternative explanation that could be considered in this context. Frequent losers might be a phenomenon found mostly in the short term: there would be a significant number of participations without victories in a short period, followed by abandonment of the public contracting market. As shown in Figure 3, presented in Section 3, companies have participated in tenders for up to ten years and were repeatedly defeated. It is not reasonable to suppose that there is a financial incentive to participate in tenders in good faith and to lose repeatedly.}

\section{It is also important to note that even if the companies participated for a short time, they still participated in a sufficient number of tenders to be classified as frequent losers by the outlier detection method.}
